%% Преамбула TeX-файла

% 1. Стиль и язык
\documentclass[utf8x, 14pt]{G7-32} % Стиль (по умолчанию будет 14pt)

\newgeometry{left=30mm, right=10mm, top=20mm, bottom=20mm}
% Остальные стандартные настройки убраны в preamble.inc.tex.
\sloppy

% Настройки стиля ГОСТ 7-32
% Для начала определяем, хотим мы или нет, чтобы рисунки и таблицы нумеровались в пределах раздела, или нам нужна сквозная нумерация.
\EqInChapter % формулы будут нумероваться в пределах раздела
\TableInChapter % таблицы будут нумероваться в пределах раздела
\PicInChapter % рисунки будут нумероваться в пределах раздела

% Добавляем гипертекстовое оглавление в PDF
\usepackage[
bookmarks=true, colorlinks=true, unicode=true,
urlcolor=black,linkcolor=black, anchorcolor=black,
citecolor=black, menucolor=black, filecolor=black,
]{hyperref}
\AfterHyperrefFix

\usepackage{microtype}% полезный пакет для микротипографии, увы под xelatex мало чего умеет, но под pdflatex хорошо улучшает читаемость

% Тире могут быть невидимы в Adobe Reader
\ifInvisibleDashes
\MakeDashesBold
\fi

\usepackage{graphicx}   % Пакет для включения рисунков

% С такими оно полями оно работает по-умолчанию:
% \RequirePackage[left=20mm,right=10mm,top=20mm,bottom=20mm,headsep=0pt,includefoot]{geometry}
% Если вас тошнит от поля в 10мм --- увеличивайте до 20-ти, ну и про переплёт не забывайте:
\geometry{right=10mm}
\geometry{left=30mm}
\geometry{bottom=20mm}
\geometry{top=20mm}

\geometry{ignorefoot}% считать от нижней границы текста

\usepackage{color} %use color
\definecolor{mygreen}{rgb}{0,0.6,0}
\definecolor{mygray}{rgb}{0.5,0.5,0.5}
\definecolor{mymauve}{rgb}{0.58,0,0.82}

\usepackage{listings} %% собственно, это и есть пакет listings
\renewcommand{\lstlistingname}{Листинг}
%%
%% GLSL definition (c) 2020 Benno Bielmeier
%%
\lstdefinelanguage{GLSL}%
{%
	morekeywords={%
	% HLSL constants
		false,FALSE,NULL,true,TRUE,%
	% GLSL predefinde macro constant
		__LINE__,__FILE__,__VERSION__,GL_core_profile,GL_es_profile,GL_compatibility_profile,%
	% GLSL precision modifier
		precision,highp,mediump,lowp,%
	% GLSL control keywords
		break,case,continue,default,discard,do,else,for,if,return,switch,while,%
	% GLSL types
		void,bool,int,uint,float,double,vec2,vec3,vec4,dvec2,dvec3,dvec4,bvec2,bvec3,bvec4,ivec2,ivec3,ivec4,uvec2,uvec3,uvec4,mat2,mat3,mat4,mat2x2,mat2x3,mat2x4,mat3x2,mat3x3,mat3x4,mat4x2,mat4x3,mat4x4,dmat2,dmat3,dmat4,dmat2x2,dmat2x3,dmat2x4,dmat3x2,dmat3x3,dmat3x4,dmat4x2,dmat4x3,dmat4x4,sampler1D,sampler2D,sampler3D,image1D,image2D,image3D,samplerCube,imageCube,sampler2DRect,image2DRect,sampler1DArray,sampler2DArray,image1DArray,image2DArray,samplerBuffer,imageBuffer,sampler2DMS,image2DMS,sampler2DMSArray,image2DMSArray,samplerCubeArray,imageCubeArray,sampler1DShadow,sampler2DShadow,sampler2DRectShadow,sampler1DArrayShadow,sampler2DArrayShadow,samplerCubeShadow,samplerCubeArrayShadow,isampler1D,isampler2D,isampler3D,iimage1D,iimage2D,iimage3D,isamplerCube,iimageCube,isampler2DRect,iimage2DRect,isampler1DArray,isampler2DArray,iimage1DArray,iimage2DArray,isamplerBuffer,iimageBuffer,isampler2DMS,iimage2DMS,isampler2DMSArray,iimage2DMSArray,isamplerCubeArray,iimageCubeArray,atomic_uint,usampler1D,usampler2D,usampler3D,uimage1D,uimage2D,uimage3D,usamplerCube,uimageCube,usampler2DRect,uimage2DRect,usampler1DArray,usampler2DArray,uimage1DArray,uimage2DArray,usamplerBuffer,uimageBuffer,usampler2DMS,uimage2DMS,usampler2DMSArray,uimage2DMSArray,usamplerCubeArray,uimageCubeArray,struct,%
	% GLSL support variables
		gl_BackColor,gl_BackLightModelProduct,gl_BackLightProduct,gl_BackMaterial,gl_BackSecondaryColor,gl_ClipDistance,gl_ClipPlane,gl_ClipVertex,gl_Color,gl_DepthRange,gl_DepthRangeParameters,gl_EyePlaneQ,gl_EyePlaneR,gl_EyePlaneS,gl_EyePlaneT,gl_Fog,gl_FogCoord,gl_FogFragCoord,gl_FogParameters,gl_FragColor,gl_FragCoord,gl_FragData,gl_FragDepth,gl_FrontColor,gl_FrontFacing,gl_FrontLightModelProduct,gl_FrontLightProduct,gl_FrontMaterial,gl_FrontSecondaryColor,gl_InstanceID,gl_Layer,gl_LightModel,gl_LightModelParameters,gl_LightModelProducts,gl_LightProducts,gl_LightSource,gl_LightSourceParameters,gl_MaterialParameters,gl_ModelViewMatrix,gl_ModelViewMatrixInverse,gl_ModelViewMatrixInverseTranspose,gl_ModelViewMatrixTranspose,gl_ModelViewProjectionMatrix,gl_ModelViewProjectionMatrixInverse,gl_ModelViewProjectionMatrixInverseTranspose,gl_ModelViewProjectionMatrixTranspose,gl_MultiTexCoord0,gl_MultiTexCoord1,gl_MultiTexCoord2,gl_MultiTexCoord3,gl_MultiTexCoord4,gl_MultiTexCoord5,gl_MultiTexCoord6,gl_MultiTexCoord7,gl_Normal,gl_NormalMatrix,gl_NormalScale,gl_ObjectPlaneQ,gl_ObjectPlaneR,gl_ObjectPlaneS,gl_ObjectPlaneT,gl_Point,gl_PointCoord,gl_PointParameters,gl_PointSize,gl_Position,gl_PrimitiveIDIn,gl_ProjectionMatrix,gl_ProjectionMatrixInverse,gl_ProjectionMatrixInverseTranspose,gl_ProjectionMatrixTranspose,gl_SecondaryColor,gl_TexCoord,gl_TextureEnvColor,gl_TextureMatrix,gl_TextureMatrixInverse,gl_TextureMatrixInverseTranspose,gl_TextureMatrixTranspose,gl_Vertex,gl_VertexID,%
	% GLSL support constants
		gl_MaxClipPlanes,gl_MaxCombinedTextureImageUnits,gl_MaxDrawBuffers,gl_MaxFragmentUniformComponents,gl_MaxLights,gl_MaxTextureCoords,gl_MaxTextureImageUnits,gl_MaxTextureUnits,gl_MaxVaryingFloats,gl_MaxVertexAttribs,gl_MaxVertexTextureImageUnits,gl_MaxVertexUniformComponents,%
	% GLSL support functions
		abs,acos,all,any,asin,atan,ceil,clamp,cos,cross,degrees,dFdx,dFdy,distance,dot,equal,exp,exp2,faceforward,floor,fract,ftransform,fwidth,greaterThan,greaterThanEqual,inversesqrt,length,lessThan,lessThanEqual,log,log2,matrixCompMult,max,min,mix,mod,noise1,noise2,noise3,noise4,normalize,not,notEqual,outerProduct,pow,radians,reflect,refract,shadow1D,shadow1DLod,shadow1DProj,shadow1DProjLod,shadow2D,shadow2DLod,shadow2DProj,shadow2DProjLod,sign,sin,smoothstep,sqrt,step,tan,texture1D,texture1DLod,texture1DProj,texture1DProjLod,texture2D,texture2DLod,texture2DProj,texture2DProjLod,texture3D,texture3DLod,texture3DProj,texture3DProjLod,textureCube,textureCubeLod,transpose,%
	% GLSL struct member -> FixMe: Should have dot(.) as delimiter
		rgb
	},
	sensitive=true,%
	comment=[l]{//},
	morecomment=[s]{/*}{*/},%
	morecomment=[l]//,%
    commentstyle=\color{purple}\ttfamily,
	stringstyle=\color{red}\ttfamily,
	morestring=[b]",%
	morestring=[b]',%
	moredelim=*[directive]\#,%
	% keyword.control.hlsl
	moredirectives={define,defined,elif,else,if,ifdef,endif,line,error,ifndef,include,pragma,undef,warning,extension,version}%
}[keywords,comments,strings,directives]%


\lstset{ %
	backgroundcolor=\color{white}, % choose the background color; you must add \usepackage{color} or \usepackage{xcolor}
	basicstyle=\footnotesize, % the size of the fonts that are used for the code
	breakatwhitespace=false, % sets if automatic breaks should only happen at whitespace
	breaklines=true, % sets automatic line breaking
	captionpos=t, % sets the caption-position to bottom
	commentstyle=\color{mygreen}, % comment style
	deletekeywords={...}, % if you want to delete keywords from the given language
	escapeinside={\%*}{*)}, % if you want to add LaTeX within your code
	extendedchars=true, % lets you use non-ASCII characters; for 8-bits encodings only, does not work with UTF-8
	frame=single, % adds a frame around the code
	keepspaces=true, % keeps spaces in text, useful for keeping indentation of code (possibly needs columns=flexible)
	keywordstyle=\color{blue}, % keyword style
	% language=Octave, % the language of the code
	morekeywords={*,...}, % if you want to add more keywords to the set
	numbers=left, % where to put the line-numbers; possible values are (none, left, right)
	numbersep=5pt, % how far the line-numbers are from the code
	numberstyle=\tiny\color{mygray}, % the style that is used for the line-numbers
	rulecolor=\color{black}, % if not set, the frame-color may be changed on line-breaks within not-black text (e.g. comments (green here))
	showspaces=false, % show spaces everywhere adding particular underscores; it overrides 'showstringspaces'
	showstringspaces=false, % underline spaces within strings only
	showtabs=false, % show tabs within strings adding particular underscores
	stepnumber=1, % the step between two line-numbers. If it's 1, each line will be numbered
	stringstyle=\color{mymauve}, % string literal style
	tabsize=2, % sets default tabsize to 2 spaces
	title=\lstname % show the filename of files included with \lstinputlisting; also try caption instead of title
}
%END of listing package%

% Пакет Tikz
\usepackage{tikz}
\usetikzlibrary{arrows,positioning,shadows}

% Произвольная нумерация списков.
\usepackage{enumerate}

% ячейки в несколько строчек
\usepackage{multirow}

% itemize внутри tabular
\usepackage{paralist,array}

%\setlength{\parskip}{1ex plus0.5ex minus0.5ex} % разрыв между абзацами
\setlength{\parskip}{1ex} % разрыв между абзацами
\usepackage{blindtext}

% Центрирование подписей к плавающим окружениям
%\usepackage[justification=centering]{caption}

\usepackage{newfloat}
\DeclareFloatingEnvironment[
placement={!ht},
name=Equation
]{eqndescNoIndent}
\edef\fixEqndesc{\noexpand\setlength{\noexpand\parindent}{\the\parindent}\noexpand\setlength{\noexpand\parskip}{\the\parskip}}
\newenvironment{eqndesc}[1][!ht]{%
    \begin{eqndescNoIndent}[#1]%
\fixEqndesc%
}
{\end{eqndescNoIndent}}


\usepackage{algorithm}
\usepackage{algpseudocode}

\newcommand{\apostrophe}{\XeTeXglyph\XeTeXcharglyph"0027\relax}



% Настройки листингов.
\ifPDFTeX
% 8 Листинги

\usepackage{listings}

% Значения по умолчанию
\lstset{
  basicstyle= \footnotesize,
  breakatwhitespace=true,% разрыв строк только на whitespacce
  breaklines=true,       % переносить длинные строки
%   captionpos=b,          % подписи снизу -- вроде не надо
  inputencoding=koi8-r,
  numbers=left,          % нумерация слева
  numberstyle=\footnotesize,
  showspaces=false,      % показывать пробелы подчеркиваниями -- идиотизм 70-х годов
  showstringspaces=false,
  showtabs=false,        % и табы тоже
  stepnumber=1,
  tabsize=4,              % кому нужны табы по 8 символов?
  frame=single
}

% Стиль для псевдокода: строчки обычно короткие, поэтому размер шрифта побольше
\lstdefinestyle{pseudocode}{
  basicstyle=\small,
  keywordstyle=\color{black}\bfseries\underbar,
  language=Pseudocode,
  numberstyle=\footnotesize,
  commentstyle=\footnotesize\it
}

% Стиль для обычного кода: маленький шрифт
\lstdefinestyle{realcode}{
  basicstyle=\scriptsize,
  numberstyle=\footnotesize
}

% Стиль для коротких кусков обычного кода: средний шрифт
\lstdefinestyle{simplecode}{
  basicstyle=\footnotesize,
  numberstyle=\footnotesize
}

% Стиль для BNF
\lstdefinestyle{grammar}{
  basicstyle=\footnotesize,
  numberstyle=\footnotesize,
  stringstyle=\bfseries\ttfamily,
  language=BNF
}

% Определим свой язык для написания псевдокодов на основе Python
\lstdefinelanguage[]{Pseudocode}[]{Python}{
  morekeywords={each,empty,wait,do},% ключевые слова добавлять сюда
  morecomment=[s]{\{}{\}},% комменты {а-ля Pascal} смотрятся нагляднее
  literate=% а сюда добавлять операторы, которые хотите отображать как мат. символы
    {->}{\ensuremath{$\rightarrow$}~}2%
    {<-}{\ensuremath{$\leftarrow$}~}2%
    {:=}{\ensuremath{$\leftarrow$}~}2%
    {<--}{\ensuremath{$\Longleftarrow$}~}2%
}[keywords,comments]

% Свой язык для задания грамматик в BNF
\lstdefinelanguage[]{BNF}[]{}{
  morekeywords={},
  morecomment=[s]{@}{@},
  morestring=[b]",%
  literate=%
    {->}{\ensuremath{$\rightarrow$}~}2%
    {*}{\ensuremath{$^*$}~}2%
    {+}{\ensuremath{$^+$}~}2%
    {|}{\ensuremath{$|$}~}2%
}[keywords,comments,strings]

% Подписи к листингам на русском языке.
\renewcommand\lstlistingname{Листинг}
\renewcommand\lstlistlistingname{Листинги}

\else
\usepackage{local-minted}
\fi


\usepackage{indentfirst}
\usepackage{ulem} % Нормальное нижнее подчеркивание

% Дополнительное окружения для подписей
\usepackage{array}

\usepackage{svg}
\usepackage{pdflscape}
\usepackage{bookmark}
% \usepackage{lscape}
% Полезные макросы листингов.
% Любимые команды
\newcommand{\Code}[1]{\textbf{#1}}

\newenvironment{signstabular}[1][1]{
	\renewcommand*{\arraystretch}{#1}
	\tabular
}{
	\endtabular
}
\floatname{algorithm}{Псевдокод}
\algrenewcommand\algorithmicwhile{\textbf{До тех пока}}
\algrenewcommand\algorithmicdo{\textbf{выполнить}}
\algrenewcommand\algorithmicrepeat{\textbf{Повторять}}
\algrenewcommand\algorithmicuntil{\textbf{Пока выполняется}}
\algrenewcommand\algorithmicend{\textbf{Конец}}
\algrenewcommand\algorithmicif{\textbf{Если}}
\algrenewcommand\algorithmicelse{\textbf{иначе}}
\algrenewcommand\algorithmicthen{\textbf{тогда}}
\algrenewcommand\algorithmicfor{\textbf{Цикл}}
\algrenewcommand\algorithmicforall{\textbf{Для всех}}
\algrenewcommand\algorithmicfunction{\textbf{Функция}}
\algrenewcommand\algorithmicprocedure{\textbf{Процедура}}
\algrenewcommand\algorithmicloop{\textbf{Зациклить}}
\algrenewcommand\algorithmicrequire{\textbf{Условия:}}
\algrenewcommand\algorithmicensure{\textbf{Обеспечивающие условия:}}
\algrenewcommand\algorithmicreturn{\textbf{Возвратить}}
\algrenewtext{EndWhile}{\textbf{Конец цикла}}
\algrenewtext{EndLoop}{\textbf{Конец зацикливания}}
\algrenewtext{EndFor}{\textbf{Конец цикла}}
\algrenewtext{EndFunction}{\textbf{Конец функции}}
\algrenewtext{EndProcedure}{\textbf{Конец процедуры}}
\algrenewtext{EndIf}{\textbf{Конец условия}}
\algrenewtext{EndFor}{\textbf{Конец цикла}}
\algrenewtext{BeginAlgorithm}{\textbf{Начало алгоритма}}
\algrenewtext{EndAlgorithm}{\textbf{Конец алгоритма}}
\algrenewtext{BeginBlock}{\textbf{Начало блока. }}
\algrenewtext{EndBlock}{\textbf{Конец блока}}
\algrenewtext{ElsIf}{\textbf{иначе если }}




\begin{document}

\frontmatter % выключает нумерацию ВСЕГО; здесь начинаются ненумерованные главы: реферат, введение, глоссарий, сокращения и прочее.


% Стиль титульного листа и заголовки


% \maketitle %создает титульную страницу


% \begin{executors}
% \personalSignature{Первый исполнитель}{ФИО}

% \personalSignature{Второй исполнитель}{ФИО}
% \end{executors}


%\listoffigures                         % Список рисунков

%\listoftables                          % Список таблиц

%\NormRefs % Нормативные ссылки 
% Команды \breakingbeforechapters и \nonbreakingbeforechapters
% управляют разрывом страницы перед главами.
% По-умолчанию страница разрывается.

% \nobreakingbeforechapters
% \breakingbeforechapters

\begin{titlepage}
    \thispagestyle{empty}

    \noindent\begin{minipage}{0.05\textwidth}
        \includegraphics[scale=0.3]{inc/img/bmstu.png}
    \end{minipage}
    \hfill
    \begin{minipage}{0.85\textwidth}\raggedleft
        \begin{center}
            \fontsize{12pt}{0.3\baselineskip}\selectfont \textbf{Министерство науки и высшего образования Российской Федерации \\ Федеральное государственное бюджетное образовательное учреждение \\ высшего образования \\ <<Московский государственный технический университет \\ имени Н.Э. Баумана \\ (национальный исследовательский университет)>> \\ (МГТУ им. Н.Э. Баумана)}
        \end{center}
    \end{minipage}

    \begin{center}
        \fontsize{12pt}{0.1\baselineskip}\selectfont
        \noindent\makebox[\linewidth]{\rule{\textwidth}{4pt}} \makebox[\linewidth]{\rule{\textwidth}{1pt}}
    \end{center}

    \begin{flushleft}
        \fontsize{12pt}{0.8\baselineskip}\selectfont

        ФАКУЛЬТЕТ \uline{
            Информатика и системы управления
            \hfill}

        КАФЕДРА \uline{\mbox{\hspace{4mm}}
            Программное обеспечение ЭВМ и информационные технологии
            \hfill}
    \end{flushleft}

    \vfill

    \begin{center}
        \fontsize{22pt}{\baselineskip}\selectfont

        \textbf{РАСЧЕТНО-ПОЯСНИТЕЛЬНАЯ ЗАПИСКА}

        \fontsize{20pt}{\baselineskip}\selectfont

        \textbf{\textit{К НАУЧНО-ИССЛЕДОВАТЕЛЬСКОЙ РАБОТЕ}}

        \textbf{\textit{НА ТЕМУ:}}
    \end{center}

    \begin{center}
        \fontsize{18pt}{0.6cm}\selectfont

        «Исследование совместимости диалектов SQL для решения задачи по конвертации запросов СУБД MySQL в PostgreSQL»
        
    \end{center}

    \vfill

    \begin{table}[h!]
        \fontsize{12pt}{0.7\baselineskip}\selectfont
        \centering
        \begin{signstabular}[0.7]{p{7.25cm} >{\centering\arraybackslash}p{4cm} >{\centering\arraybackslash}p{4cm}}
            Студент группы ИУ7-56Б & \uline{\mbox{\hspace*{4cm}}} & \uline{\hfill Д.В. Варин  \hfill} \\
            & \scriptsize (Подпись, дата) & \scriptsize (И.О. Фамилия)
        \end{signstabular}

        \vspace{\baselineskip}

        \begin{signstabular}[0.7]{p{7.25cm} >{\centering\arraybackslash}p{4cm} >{\centering\arraybackslash}p{4cm}}
            Руководитель & \uline{\mbox{\hspace*{4cm}}} & \uline{\hfill Д.А. Кузнецов \hfill} \\
            & \scriptsize (Подпись, дата) & \scriptsize (И.О. Фамилия)
        \end{signstabular}


        % \begin{signstabular}[0.7]{p{7.25cm} >{\centering\arraybackslash}p{4cm} >{\centering\arraybackslash}p{4cm}}
        %     Консультант & \uline{\mbox{\hspace*{4cm}}} & \uline{\hfill  \hfill} \\
        %     & \scriptsize (Подпись, дата) & \scriptsize (И.О. Фамилия)
        % \end{signstabular}
    \end{table}

    \vfill

    \begin{center}
        \normalsize \textit{\textbf{2022} г.}
    \end{center}
\end{titlepage}

% Также можно использовать \Referat, как в оригинале
\begin{abstract}

    Отчет содержит \pageref{LastPage}\,стр.%
    \ifnum \totfig >0
    , \totfig~рис.%
    \fi
    \ifnum \tottab >0
    , \tottab~табл.%
    \fi
    %
    \ifnum \totbib >0
    , \totbib~источн.%
    \fi
    %
    \ifnum \totapp >0
    , \totapp~прил.%
    \else
    %
    \fi
    \nocite{*}

\end{abstract}

%%% Local Variables: 
%%% mode: latex
%%% TeX-master: "rpz"
%%% End: 


\tableofcontents

\printnomenclature % Автоматический список сокращений

\Introduction
Реляционная модель данных нужна для организации данных в виде таблиц (отношений) с помощью связей - полей таблиц, 
ссылающихся на другие таблицы, которые называются внешними ключами.  
Сегодня существуют другие модели данных, включая NoSQL, но системы управления реляционными базами данных (СУБД) остаются доминирующими для хранения 
и управления данными во всем мире \cite{db-engines-rating}. 

Одними из самых популярных СУБД являются MySQL и PostgreSQL. 

MySQL --- свободная реляционная система управления базами данных. 

PostgreSQL - свободная объектно-реляционная система управления базами данных.
Postgres более функциональная СУБД, нежели MySQL хотя и сложнее для использования.


Часто, с развитием проекта, возникает необходимость миграции с одного движка СУБД на другой. 
Одной из причин является нехватка функциональных возможностей.  

PostgreSQL является популярным выбором для функций NoSQL. 
Она изначально поддерживает большое разнообразие типов данных, включая JSON, hstore и XML \cite{postgres-types}.

Управление параллельным доступом посредством многоверсионности (MVCC) -  это одна из главных причин, почему компании выбирают PostgreSQL. 
MVCC предоставляет одновременный доступ к базе данных множеству агентов на чтение и запись. 
Это устраняет необходимость каждый раз блокировать чтение-запись, когда кто-то взаимодействует с данными. 
Таким образом, значительно повышается эффективность управления СУБД и ее производительность.

MVCC обеспечивает такую функциональность через изоляцию снапшотов. 
Моментальные снимки (снапшоты) представляют состояние данных в определенный момент времени.

Цель работы --- исследовать совместимость диалектов sql для СУБД MySQL и PostgreSQL для решения задачи по конвертации запросов.
Таким образом, необходимо выполнить следующие задачи:
\begin{itemize}
\item выяснить совместимость СУБД с стандартом SQL;
\item выявить различия диалектов PostgreSQL, MySQL;
\item определить взаимозаменяемость конструкций;
\item описать существующие решения по конвертации запросов MySQL в PostgreSQL;
\item выяснить особенности и их преимущества;
\item сформулировать выводы.
\end{itemize}

\Abbrev{СУБД}{Система управления базами данных "" --- совокупность программных средств, обеспечивающих управление созданием и использованием баз данных}
\Abbrev{SQL}{Structured Query Language ""--- язык структурированных запросов}
\Abbrev{NoSQL}{Not only Structed Query Language""--- термин, обозначающий ряд подходов, направленных на реализацию хранилищ баз данных, имеющих существенные отличия от моделей, используемых в традиционных реляционных СУБД с доступом к данным средствами языка SQL}
\Abbrev{MVCC}{Multiversion Concurrency Control ""--- один из механизмов СУБД для обеспечения параллельного доступа к базам данных, заключающийся в предоставлении каждому пользователю так называемого «снимка» базы, обладающего тем свойством, что вносимые пользователем изменения невидимы другим пользователям до момента фиксации транзакции}
\Abbrev{JSON}{JavaScript Object Notation ""--- текстовый формат обмена данными, основанный на JavaScript}
\Abbrev{XML}{Extensible Markup Language ""--- расширяемый язык разметки, используемый для хранения и передачи данных}
\Abbrev{ANSI}{American National Standards Institute ""--- объединение американских промышленных и деловых групп, разрабатывающее торговые и коммуникационные стандарты}
\Abbrev{ISO}{International Organization for Standardization ""--- международная организация, занимающаяся выпуском стандартов}



\Define{PostgreSQL}{свободная объектно-реляционная система управления базами данных}
\Define{MySQL}{свободная реляционная система управления базами данных}
\Define{Microsoft SQL Server}{cистема управления реляционными базами данных, разработанная корпорацией Microsoft}







\mainmatter % это включает нумерацию глав и секций в документе ниже

\chapter{Анализ предметной области}
\label{cha:body}
Базы данных --- это упорядоченный набор структурированной информации или данных, которые обычно хранятся в электронном виде в компьютерной системе. 
База данных обычно управляется системой управления базами данных (СУБД). 
Система управления базами данных --- это программа, которая взаимодействует с базой данных \cite{oracle-db}. 
СУБД позволяет контролировать доступ к базе данных, записывать данные, выполнять запросы и любые другие задачи, связанные с управлением базой данных. 

База данных может быть любым набором данных, а не только хранящимся на компьютере, СУБД --- это программное обеспечение, которое позволяет взаимодействовать с базой данных.

Большинство реляционных баз данных используют язык структурированных запросов (SQL) для управления данными и запросов к ним.
Однако многие СУБД используют свой собственный диалект SQL, который может иметь определенные ограничения или расширения. 
Эти расширения обычно включают в себя дополнительные функции, которые позволяют пользователям выполнять более сложные операции, 
чем они могли бы делать, использую стандартный SQL. 

Стандартный SQL --- язык структурированных запросов, поддерживаемый организациями ANSI\cite{ansi-sql} и ISO\cite{iso-sql}. 
В стандарте используется модульная структура, основная часть стандарта вынесена в раздел "SQL/Foundation".  
Для современных СУБД существует требование по соблюдению стандарта --- они должны поддерживать основной модуль. 
Однако полной поддержки стандарта не предоставляет ни одна из систем.  
Так, стандарт SQL 2016 года содержит 177 функций в основном модуле. PostgreSQL реализует 170 из них \cite{postgresql-standard},
MySQL --- чуть более 140 \cite{mysql-standard}.  

По статистике на январь 2022 года \cite{sql-rating}, MySQL занимает 2 место, Microsoft SQL Server 3 место в мире по популярности 
среди реляционных СУБД. PostgreSQL находится на 4 месте. 


Исходя из больших функциональных возможностей и лучшей поддержкой управления параллельным доступом, 
переход с MySQL на PostgreSQL является популярной задачей разработки. 

В научно-исследовательской работе предстоит выяснить, чем отличаются данные СУБД, совместимы ли они, 
как перейти с одной на другую.  

% %%% Local Variables:
% %%% mode: latex
% %%% TeX-master: "rpz"
% %%% End:


\chapter{Различия в диалектах SQL СУБД MySQL и PostgreSQL}
\label{cha:classification}
Для полноценной конвертации запросов следует выявить особенности, которые есть в MySQL, но нет в PostgreSQL,
различия в данных СУБД и привести способ адаптирования запросов MySQL в PostgreSQL.
Ниже приведены сравнительные таблицы на основе данных, полученных в ходе знакомства с документацией СУБД, а также 
некоторых статей \cite{rdbms-standard-diff}, \cite{cmp-postgresql-mysql}.
\section{Функциональные различия}
Ниже приведена сравнительная таблица функциональных различий между СУБД.  
В первую очередь учитываются моменты, которые есть в MySQL, но нет в PostgreSQL. При конвератции запросов необходимо обратить особое внимание
на описанные ниже различия и адаптировать их под новую систему.
\clearpage
    \begin{table}[]
        \small
        \caption{Функциональные различия СУБД}
        \begin{tabular}{|l|l|l|}
        \hline
        Возможности                                                                 & MySQL                                                                                                           & PostgreSQL                                                                                                   \\ \hline
        Работа с датой                                                              & \begin{tabular}[c]{@{}l@{}}CURDATE() \\ CURTIME() \\ EXTRACT()\end{tabular}                                     & \begin{tabular}[c]{@{}l@{}}CURRENT\_DATE() \\ CURRENT\_TIME() \\ EXTRACT()\end{tabular}                      \\ \hline
        Autoincrement тип                                                           & auto\_increment                                                                                                 & serial                                                                                                       \\ \hline
        Столбец идентификации                                                       & Нет                                                                                                             & Да                                                                                                           \\ \hline
        \begin{tabular}[c]{@{}l@{}}Аналитические \\ функции\end{tabular}            & Нет                                                                                                             & Да                                                                                                           \\ \hline
        \begin{tabular}[c]{@{}l@{}}Установка параметра \\ по умолчанию\end{tabular} & Константа                                                                                                       & Константа и функция                                                                                          \\ \hline
        \begin{tabular}[c]{@{}l@{}}ЯП для хранимых \\ процедур\end{tabular}         & \begin{tabular}[c]{@{}l@{}}Синтаксис, \\ удовлетворяющий SQL:2003\end{tabular}                                  & \begin{tabular}[c]{@{}l@{}}Ruby, Perl, Python, TCL, \\ PL/pgSQL, \\ SQL, JavaScript и др.\end{tabular}       \\ \hline
        Индексы                                                                     & \begin{tabular}[c]{@{}l@{}}Только индексы:\\ B-Tree, R-Tree, Hash, Inverted\end{tabular}                        & \begin{tabular}[c]{@{}l@{}}Индексы MySQL + \\ Частичные, Bitmap, Expression, \\ Функциональные\end{tabular}  \\ \hline
        Условные конструкции                                                        & \begin{tabular}[c]{@{}l@{}}IF, \\ IFNULL\end{tabular}                                                           & CASE                                                                                                         \\ \hline
        \begin{tabular}[c]{@{}l@{}}Возможность удалить \\ TEMP TABLE\end{tabular}   & Да                                                                                                              & Нет                                                                                                          \\ \hline
        \end{tabular}
        \end{table}
Во время поиска функциональных различий MySQL и PostgreSQL были получен следующие выводы:  
конвертация запросов с MySQL на PostgreSQL проще в реализации, чем наоборот, так
как MySQL реализует меньшее количество стандартных SQL функций, чем PostgreSQL.
Это накладывает трудности при конвертации запросов, следует
выявить конкретные различия в используемых запросах и написать аналоги для PostgreSQL. 
MySQL не поддерживает следующие возможности:
\begin{itemize}
    \item объединение таблиц FULL OUTER JOINS;
    \item пересечение - INTERSECT;
    \item разность - EXCEPT;
    \item невозможно выполнить операторы UPDATE или DELETE с SELECT подзапросом из той же таблицы;
    \item ограничение количества строк LIMIT не поддерживается в подзапросах;
    \item подзапрос LIMIT с IN/ALL/ANY/SOME;
    \item частичные индексы, битовые индексы и индексы выражений;
    \item оконные функции нельзя использовать как часть UPDATE или DELETE;
    \item ограничение на уникальность строк DISTINCT не поддерживается функциями окна;
    \item вложенные оконные функции не поддерживаются.
\end{itemize} 

\section{Синтаксические различия}
Ниже приведена таблица синтаксических различий СУБД.
\begin{table}[]
    \small
    \caption{Синтаксические различия СУБД}
        \begin{tabular}{|l|l|l|}
        \hline
        MySQL                                                                  & PostgreSQL                                                                         & Комментарий                                                                                                                                                                                                                                                                                                                                                                                                                    \\ \hline
        \#                                                                     & --                                                                                 & \begin{tabular}[c]{@{}l@{}}MySQL использует \# для начала строки комментария;\\PostgreSQL использует двойное тире (стандарт ANSI);\\ MySQL может использовать --, но требует пробела\\ после --, \\ в то время как в PostgreSQL это не обязательно.\end{tabular}                                                                                                                                                                \\ \hline
        \begin{tabular}[c]{@{}l@{}} \apostrophe Vasiliy\apostrophe\ и \\ \apostrophe \apostrophe Vasiliy\apostrophe \apostrophe \end{tabular}        & \apostrophe Vasiliy\apostrophe & \begin{tabular}[c]{@{}l@{}}MySQL использует \apostrophe или \apostrophe \apostrophe\ для заключения значений\\в кавычки.\\Это не стандарт ANSI для баз данных. PostgreSQL\\ использует для этого только одинарные кавычки.\\ Двойные кавычки используются для обозначения\\ системных идентификаторов, имен полей\\имен таблиц и т. д.. С помощью SET\\ sql\_mode=\apostrophe ANSI\_QUOTES\apostrophe\ можно заставить MySQL \\ использовать только одинарные кавычки\end{tabular} \\ \hline
        \begin{tabular}[c]{@{}l@{}}... WHERE \\ lastname=\\ \apostrophe \apostrophe car\apostrophe \apostrophe\ \end{tabular} & \begin{tabular}[c]{@{}l@{}}... WHERE \\ lower(lastname)=\\ \apostrophe car\apostrophe\end{tabular}      & \begin{tabular}[c]{@{}l@{}}PostgreSQL чувствителен к регистру\\ при сравнении строк. MySQL - нет.\end{tabular}                                                                                                                                                                                                                                                                                                                 \\ \hline
        \begin{tabular}[c]{@{}l@{}}\apostrophe CaR\apostrophe\ = \apostrophe car\apostrophe\\ а может и нет\end{tabular}  & \apostrophe \apostrophe CaR\apostrophe \apostrophe\ \textless{}\textgreater \apostrophe \apostrophe car\apostrophe \apostrophe & \begin{tabular}[c]{@{}l@{}}Имена баз данных, таблиц, полей и столбцов\\ в PostgreSQL не зависят от регистра, если их не создали\\ с двойными кавычками вокруг имени.\\ В этом случае они чувствительны к регистру.\\ \\ В MySQL имена таблиц могут быть чувствительны \\ к регистру или нет,в зависимости от \\ используемой операционной системы.\end{tabular}                                                               \\ \hline
        \begin{tabular}[c]{@{}l@{}}\apostrophe foo\apostrophe\ || \apostrophe bar\apostrophe\ \\ означает или\end{tabular} & \begin{tabular}[c]{@{}l@{}}\apostrophe foo\apostrophe\ || \apostrophe bar\apostrophe\ \\ означает \\ конкатенацию\end{tabular} & \begin{tabular}[c]{@{}l@{}}MySQL принимает операторы языка C для логики, \\ SQL требует AND, OR. Решение - использовать\\ стандартные ключевые слова SQL,\\ обе базы данных понимают это.\end{tabular}                                                                                                                                                                                                                         \\ \hline
        \end{tabular}
    \end{table}
\clearpage
\section{Различия типов данных}
Ниже приведена таблица сопоставления типов данных MySQL и PostgreSQL.
\begin{table}[]
    \small
    \caption{Таблица сопоставления типов данных СУБД}
        \begin{tabular}{|l|l|l|}
        \hline
        MySQL                                                                                                                              & PostgreSQL                                                                                                                              & Комментарий                                                                                                                                                                      \\ \hline
        \begin{tabular}[c]{@{}l@{}}TINY(SMALL)INT\\ MEDIUMINT\\ BIGINT\end{tabular}                                                    & \begin{tabular}[c]{@{}l@{}}SMALLINT\\ INTEGER\\ BIGINT\end{tabular}                                                          & \begin{tabular}[c]{@{}l@{}}Размер INTEGER в PostgreSQL\\ 4 байта\end{tabular}                                                                                                    \\ \hline
        \begin{tabular}[c]{@{}l@{}}TINYINT UNSIGNED\\ SMALLINT UNSIGNED\\ MEDIUMINT UNSIGNED\\ INT UNSIGNED\\ BIGINT UNSIGNED\end{tabular} & \begin{tabular}[c]{@{}l@{}}SMALLINT\\ INTEGER\\ INTEGER\\ BIGINT\\ NUMERIC(20)\end{tabular}                                             & \begin{tabular}[c]{@{}l@{}}PostgreSQL соответствует\\ стандарту, в нем нет \\ беззановых типов\end{tabular}                                                                      \\ \hline
        \begin{tabular}[c]{@{}l@{}}FLOAT(UNSIGNED)\end{tabular}                                                                     & \begin{tabular}[c]{@{}l@{}}REAL\end{tabular}                                                                                     & \begin{tabular}[c]{@{}l@{}}PostgreSQL чувствителен к регистру \\ при сравнении строк. MySQL - нет.\end{tabular}                                                                  \\ \hline
        DOUBLE                                                                                                                             & DOUBLE PRECISION                                                                                                                        &                                                                                                                                                                                  \\ \hline
        BOOLEAN                                                                                                                            & BOOLEAN                                                                                                                                 & \begin{tabular}[c]{@{}l@{}}В MySQL тип BOOLEAN - это \\ псевдоним для TINYINT(1).\\ PostgreSQL не поддерживает \\ автоматическую\\ конвертацию bool в числовые типы\end{tabular} \\ \hline
        \begin{tabular}[c]{@{}l@{}}(TINY,MEDIUM,\\LONG)TEXT\end{tabular}                                                    & \begin{tabular}[c]{@{}l@{}}TEXT\end{tabular}                                                                       &                                                                                                                                                                                  \\ \hline
        \begin{tabular}[c]{@{}l@{}}(VAR)BINARY(n)\\ (TINY)BLOB\\ (MEDIUM,LONG)BLOB\end{tabular}                         & \begin{tabular}[c]{@{}l@{}}BYTEA\\ BYTEA\\ BYTEA\\ BYTEA\end{tabular}                                                   &                                                                                                                                                                                  \\ \hline
        ZEROFILL                                                                                                                           & Не поддерживается                                                                                                                       &                                                                                                                                                                                  \\ \hline
        \begin{tabular}[c]{@{}l@{}}DATE\\ TIME\\ DATETIME\\ TIMESTAMP\end{tabular}                                                         & \begin{tabular}[c]{@{}l@{}}DATE\\ TIME {[}NO TZ{]}\\ TIMESTAMP {[}NO TZ{]}\\ TIMESTAMP {[}NO TZ{]}\end{tabular} &                                                                                                                                                                                  \\ \hline
        \end{tabular}
    \end{table}
\clearpage
\section{Анализ различий для решения задачи конвертации запросов}

При рассмотрении способов конвертации запросов, следует рассмотреть два варианта - с ограничениями и без них. 
Это необходимо, потому что каждая СУБД имеет свои уникальные функции, которых нет в других СУБД или их функционал не в точности повторяется.

\subsection{Конвертация с ограничениями}
Существующие различия между СУБД препятствуют реализации полной конвертации запросов.

Если бы стандарт SQL был реализован полностью каждой из рассматриваемых СУБД, или же были реализованы
одни и те же функции, то одним из способов конвертации запросов(при условии, что СУБД 
использует только принятые стандартом SQL функций) был бы следующий:  
так как все используемые функции являются стандартными, то следует перевести синтаксис MySQL на PostgreSQL,
тем самым решив задачу конвертации.

Однако использование каждой СУБД исключительно функций из стандарта SQL убирает их преимущества, 
отличающие каждую СУБД.

\subsection{Полная конвертация запросов}
Для решения задачи конвертации запросов, следует выполнить следующие действия:
\begin{itemize}
    \item заменить несовместимые типы данных на аналоги в PostgreSQL;
    \item перевести запросы с синтаксиса MySQL на синтаксис PostgreSQL;
    \item заменить нестандартные функции на аналоги в PostgreSQL;
    \item СУБД MySQL нечувствительна к регистру, следует определить общий регистр;
    \item адаптировать отличающиеся типы.
\end{itemize} 

В результате запросы для MySQL будут преобразованы в корректные запросы для PostgreSQL и задача конвертации будет решена. 



\backmatter %% Здесь заканчивается нумерованная часть документа и начинаются ссылки и
            
\Conclusion % заключение к отчёту
 
В научно-исследовательской работе были проанализированы различия и на их основе составлен перечень, 
на основе которого можно конвертировать запрос из СУБД MySQL в PostgreSQL.

Было выявлено, что диалекты несовместимы, так как они имеют отличающиеся функциональные возможности, которые
реализованы сверх стандарта SQL. При этом PostgreSQL предоставляет больше возможностей по взаимодействию с данными 
и имеет более полный функционал.  
Также в результате исследования удалось выяснить различия между СУБД, определить
взаимозаменяемые конструкции и привести сравнительные таблицы,

В результате работы было сформулировано решение для конвертации запроса из MySQL в PostgreSQL 
с использованием полученной в работе информации. 

%%% Local Variables:
%%% mode: latex
%%% TeX-master: "rpz"
%%% End:
 
 

%% заключение


% % Список литературы при помощи BibTeX
% Юзать так:
%
% pdflatex rpz
% bibtex rpz
% pdflatex rpz

\bibliographystyle{ugost2008}
\bibliography{rpz}

%%% Local Variables: 
%%% mode: latex
%%% TeX-master: "rpz"
%%% End: 



% \appendix   % Тут идут приложения

% \include{90-appendix1}

% \include{91-appendix2}

\end{document}

%%% Local Variables:
%%% mode: latex
%%% TeX-master: t
%%% End:
